\chapter{Introducción al Trabajo Fin de Máster}
% BORRADOR para editar a tu voz (estilo impersonal UAX). Recomendado: 5-10 páginas.

La interpretación de imágenes médicas es una de las tareas de mayor valor y, a la vez, de mayor
carga asistencial dentro del sistema sanitario. La radiografía de tórax (CXR) es la prueba de
imagen más solicitada del mundo, y la redacción de su informe consume una parte considerable del
tiempo del radiólogo. En este contexto, la generación automática de informes mediante inteligencia
artificial se presenta como una vía prometedora de asistencia. Sin embargo, para que un sistema de
este tipo sea aceptable en la práctica clínica no basta con que produzca un texto plausible: debe
ser posible saber \textit{cuándo fiarse} de él y \textit{por qué} concluye lo que concluye. Este
capítulo introduce la motivación del trabajo, el problema que aborda y la organización de la
memoria.

\section{Justificación del proyecto realizado}

Mi pareja se encontraba realizando un grado de Imagen para el Diagnóstico y Medicina Nuclear, y debido
a esto tuve la oportunidad de hablar con algunos de los profesores del centro sobre cómo la tecnología
y, en concreto, la inteligencia artificial estaban teniendo un impacto importante en la profesión. De
ahí, la conversación derivó a los principales problemas que tenían con ella (sobre todo la generativa,
no la clasificadora clásica, que lleva muchísimo tiempo en la industria y va por una rama completamente
distinta a la que usa VLMs), y el tema que más se repetía era la falta de explicabilidad, lo ``caja
negra'' que es, que impide a los radiólogos comprender el porqué del diagnóstico, ya que, a diferencia
de los modelos clasificadores, que simplemente clasifican una patología (de manera correcta o
incorrecta), los VLMs dan una explicación de lo que ven en la imagen para apoyar el diagnóstico, pero,
muchas veces, esto resulta en alucinaciones, que producen el efecto contrario al que se quiere generar:
en este caso, agilizar y facilitar el diagnóstico mediante estas radiografías, que, al tener estas
``alucinaciones'', o errores de interpretación, hacen que se dedique más tiempo a la revisión. Esto hizo
que mi planteamiento inicial fuese intentar encontrar una manera de ``explicar'' el diagnóstico del
modelo mediante mapas de calor que interpretasen la ``atención'' del modelo hacia distintas zonas de la
imagen. Si bien, por las propias circunstancias del proyecto, se acabó derivando gran parte del esfuerzo
a otras áreas, la intención seguía siendo la misma: facilitar la comprensión y aumentar la confianza de
los profesionales sanitarios en los VLMs mediante explicabilidad.

Desde el punto de vista de su utilidad, el proyecto se sitúa en la intersección de tres
necesidades actuales. En el ámbito \textbf{clínico}, la presión asistencial y la escasez de
radiólogos hacen atractiva cualquier herramienta que agilice el flujo de trabajo sin comprometer
la seguridad del paciente. En el ámbito \textbf{científico-técnico}, los modelos multimodales de
lenguaje (MLLM) han alcanzado una notable competencia en tareas de imagen y texto, pero su
aplicación a dominios de alta responsabilidad exige garantías de fiabilidad y trazabilidad de las
que todavía carecen. Y en el ámbito \textbf{ético y regulatorio}, la interpretabilidad de los
sistemas de IA en salud no es una característica opcional sino un requisito: un modelo cuyas
decisiones no pueden auditarse difícilmente será adoptado. El presente trabajo se orienta a este
espacio, estudiando no solo la generación de informes sino, sobre todo, la \textbf{fidelidad y el
anclaje clínico de sus explicaciones} y la posibilidad de estimar la \textbf{confianza} en cada
informe generado.

\section{Presentación del proyecto y sus contenidos}
El problema que se aborda puede enunciarse de forma simple: un sistema que genera informes
radiológicos produce, para cada afirmación, una explicación visual (un mapa de calor que señala la
región de la imagen que la motivó); la cuestión es si esa explicación es \textit{fiel} al
razonamiento del modelo, si está \textit{anclada} a la anatomía correcta y, en última instancia,
si permite \textbf{detectar cuándo el modelo se equivoca}.

Considérese, a modo de ilustración, un sistema que informa ``cardiomegalia'' sobre una radiografía
de tórax normal. Si junto al texto se muestra un mapa de calor razonablemente centrado sobre la
silueta cardíaca, el radiólogo dispone de una pista de \emph{dónde} ha mirado el modelo, pero no de
ninguna garantía de que ese mismo patrón de atención no acompañe también a sus aciertos: una
explicación visualmente plausible puede convivir con un contenido incorrecto. La pregunta que
articula el trabajo es, precisamente, si existe alguna señal —en la propia explicación o en la
confianza interna del modelo— capaz de distinguir este caso del de un informe verdaderamente
correcto, antes de que el error llegue al informe final.

Para responderla, el trabajo contribuye en cuatro frentes. En primer lugar, se desarrolla un
\textbf{generador propio} encoder--decoder multimodal, en el que la restricción de datos se
convierte en una decisión metodológica mediante la destilación de conocimiento desde modelos
preentrenados libres. En segundo lugar, se construye un \textbf{marco de explicabilidad} que
evalúa dos ejes: la fidelidad de las explicaciones y su anclaje clínico. En tercer lugar, dado que
un modelo de escala reducida limita el alcance del análisis, se \textbf{eleva el estudio a un MLLM
del estado del arte} (LLaVA-Rad) y se valida su anclaje de forma \textbf{objetiva} contra
anotaciones reales de localización. Por último, se investiga si la explicación —o la propia
incertidumbre del modelo— permite \textbf{detectar alucinaciones}, y se comprueba que la señal de
confianza resultante es \textbf{agnóstica al modelo y robusta al dominio} mediante una réplica en un
segundo generador (MAIRA-2) y un segundo dominio (PadChest-GR); la señal se \textbf{calibra} como una
confianza por informe y se materializa en una demostración interactiva que integra informe, anclaje y
fiabilidad.

El interés del problema se ha intensificado con la irrupción de los \textbf{modelos multimodales de
lenguaje} (MLLM), que acoplan un potente codificador visual a un modelo de lenguaje de gran tamaño y
han elevado de forma notable la calidad del texto generado a partir de imágenes. Su aplicación a la
radiografía de tórax promete asistir en la redacción del informe, agilizando el flujo de trabajo. Sin
embargo, esta misma potencia trae consigo un riesgo característico: la \textbf{alucinación}, es decir,
la generación de afirmaciones plausibles pero no sustentadas por la imagen. En un dominio de alta
responsabilidad como el diagnóstico por imagen, una alucinación no detectada puede tener consecuencias
graves, y una omisión —no reportar un hallazgo presente— otras distintas pero igualmente serias. Por
ello, el problema no es solo generar buen texto, sino \textbf{saber cuándo el texto generado es
fiable}, y hacerlo de una forma que el clínico pueda interpretar y en la que pueda confiar. Es en esta
intersección —entre la capacidad generativa de los MLLM y la exigencia de fiabilidad y transparencia
del entorno clínico— donde se inscribe el presente trabajo.

\section{Preguntas de investigación}
El problema anterior se concreta en cuatro preguntas de investigación que vertebran el trabajo:
\begin{itemize}
  \item \textbf{PI1.} ¿Está la atención de un MLLM de radiografía de tórax \emph{anclada} sobre la
  región anatómica de la lesión, contrastada con anotaciones objetivas de localización?
  \item \textbf{PI2.} ¿Permite la \emph{morfología} de la explicación (la forma del mapa de calor, o
  las cajas de anclaje del propio modelo) detectar cuándo el modelo se equivoca?
  \item \textbf{PI3.} ¿Permite la \emph{incertidumbre} de la decisión del modelo detectar sus errores,
  y con qué asimetría entre alucinaciones y omisiones?
  \item \textbf{PI4.} ¿Es esa señal de fiabilidad \emph{agnóstica al modelo y robusta al dominio}, es
  decir, se transfiere a otro generador y a otra población?
\end{itemize}

\clearpage
\section{Contribuciones}
Las principales contribuciones del trabajo son las siguientes. Primera, una \textbf{auditoría
objetiva del anclaje} de la atención de un MLLM de CXR frente a cajas delimitadoras reales. Segunda,
la constatación —por triple vía, incluida la de las cajas entrenadas del propio modelo— de que la
\textbf{morfología de la explicación no predice la fiabilidad}, así como la identificación de varios
\emph{artefactos} metodológicos que solo un control cuidadoso de confusores reveló. Tercera, la
\textbf{disociación} entre errores detectables por incertidumbre (alucinaciones) y errores confiados
(omisiones), y la propuesta de una señal de confianza \textbf{simple, agnóstica y disponible en
inferencia}. Cuarta, una \textbf{validación cruzada} de dicha señal en un diseño de dos generadores
por dos dominios, que demuestra su robustez. Y quinta, su \textbf{calibración} y su traducción en una
demostración interactiva que integra informe, anclaje y confianza para el clínico. Como aportación
metodológica menor, se documenta que \emph{excluir el token de apertura del informe} mejora de forma
sistemática la señal de incertidumbre.

La memoria se organiza como sigue. El Capítulo~2 fija el objetivo general y los específicos, y
explica cómo evolucionó el alcance del proyecto hasta llegar a ellos. El Capítulo~3 revisa el estado
del arte en cuatro frentes: generación automática de informes, funcionamiento interno de los modelos
multimodales, explicabilidad de \textit{transformers} y detección de incertidumbre y alucinaciones. El
Capítulo~4 describe en detalle la metodología: la población y la muestra, las hipótesis, las variables
y los instrumentos de medida, la arquitectura del generador propio y el diseño experimental. El
Capítulo~5 presenta los resultados siguiendo el orden en que se han enunciado las preguntas de
investigación, y los discute a la luz de la literatura y de sus implicaciones prácticas. El
Capítulo~6, por último, recoge las conclusiones, valora las limitaciones —incluidas las derivadas de
las condiciones materiales del proyecto— y traza las líneas futuras, tanto de investigación como de
mejora del sistema.
